% ch5.tex -- Module 8, Chapter 5: Deep inside the atom (Standard Model, particle accelerators)
% All figures are TikZ / pgfplots, drawn inline.
\chapter[Deep Inside the Atom]{Deep Inside the Atom: the Standard Model and Particle Accelerators}\label{ch5:chap}
\covers{M8.IQ5.01, M8.IQ5.02, M8.IQ5.03}

% chapter-local macros (prefix chFive, as the style guide asks)
\newcommand\chFiveaq[1]{\ensuremath{\overline{\mathrm{#1}}}} % antiquark, e.g. \chFiveaq{d}
\newcommand\chFiveperc{\ensuremath{/c^2}}                     % "GeV/c^2" after \qty{..}{GeV}
\newcommand\chFiveanu{\ensuremath{\overline{\nu}_\mathrm{e}}} % electron antineutrino

\inquiry{How is it known that human understanding of matter is still incomplete?}

\begin{goals}
  \item analyse the evidence that protons and neutrons are made of quarks, and that
        other subatomic particles exist
  \item describe the Standard Model: quarks, leptons, hadrons and the four forces
  \item explain how linacs, cyclotrons and synchrotrons work, using $W = qV$,
        $r = mv/(qB)$ and $T = 2\pi m/(qB)$
  \item assess how accelerators test the Standard Model, and why our understanding of
        matter is still incomplete
\end{goals}

%% ==================================================================
\section{Are protons and neutrons fundamental?}\label{ch5:sec:evidence}
\covers{M8.IQ5.01}

\begin{bigidea}
By 1932 (when the neutron was found), every atom could be explained with protons, neutrons
and electrons. Later experiments showed that protons and neutrons are made of smaller
particles, called quarks. They also found many other subatomic particles (particles
smaller than an atom).
\end{bigidea}

\begin{keywords}
  \kw{Fundamental particle}{a particle with no smaller parts, as far as we can tell.}
  \kw{Deep inelastic scattering}{firing very fast electrons at protons to see inside them.}
  \kw{Particle zoo}{the hundreds of new particles found in the 1950s and 1960s.}
  \kw{Antiparticle}{a partner with the same mass but opposite charge and properties.}
\end{keywords}

A \term{fundamental particle} is a particle with no smaller parts. The electron is
fundamental: experiments find no size and no parts inside it. A proton or a neutron is
different. Three kinds of evidence show that they have parts.

\subsection{Evidence 1: deep inelastic scattering}

From 1967 to 1973, physicists at SLAC (Stanford, USA) fired electrons with energies up
to about \qty{20}{GeV} at protons. This is \term{deep inelastic scattering}. ``Deep''
means the electrons look deep inside the proton. ``Inelastic'' means the electrons lose
a lot of kinetic energy, and the proton breaks apart. It is Rutherford's gold-foil idea
(\cref{ch2:sec:nucleus}), one level deeper.

\Cref{ch5:fig:dis} compares two models. Suppose the proton's charge were smeared evenly
through it, like jam soaked into a sponge. Then each electron would get only a small
sideways pull. But SLAC found something else. Far more electrons than expected came out
at large angles. So the proton contains tiny, \term{point-like}
charges: charges so small that they act like dots with no size. We now call them
\term{quarks}: the particles that make up protons and neutrons.

\begin{figure}
  \centering
  \begin{tikzpicture}[x=0.85cm, y=0.85cm,
      etrack/.style={line width=0.9pt, color=hscvel, midarrow},
      eout/.style={line width=0.9pt, color=hscvel, -{Stealth[length=2.2mm, width=1.7mm]}},
      qk/.style={circle, text=white, font=\tiny\bfseries, inner sep=0pt,
        minimum size=0.36cm}]
    % Model A: charge spread evenly. Electrons (negative) are pulled slightly TOWARDS the
    % positive charge, so every track bends towards the axis.
    \node[font=\small\bfseries, text=hscink] at (0,1.7) {Model A: smooth proton};
    \fill[hscforce!14] (0,0) circle (1.05);
    \draw[hscforce!70, line width=0.7pt] (0,0) circle (1.05);
    \draw[etrack] (-2.4,0.75) -- (-0.4,0.75) -- (2.4,0.55);
    \draw[etrack] (-2.4,0.35) -- (-0.2,0.35) -- (2.4,0.24);
    \draw[etrack] (-2.4,-0.45) -- (-0.2,-0.45) -- (2.4,-0.31);
    \draw[etrack] (-2.4,-0.85) -- (-0.4,-0.85) -- (2.4,-0.64);
    \node[font=\scriptsize, text=hscforce!80!black, fill=hscforce!14, inner sep=1pt]
      at (0,-0.05) {$+$ charge spread evenly};
    \node[font=\footnotesize, text=hscmuted] at (0,-1.75)
      {prediction: only small deflections};
    \node[diagramlabel, text=hscvel, anchor=east] at (-2.45,0) {electrons};
    % Model B: point-like charges. One electron is turned back by about 150 degrees
    % (by the negative d), one leaves at a large angle (u).
    \begin{scope}[xshift=6.4cm]
      \node[font=\small\bfseries, text=hscink] at (0,1.7) {Model B: tiny charges inside};
      \draw[hscmuted, dashed, line width=0.7pt] (0,0) circle (1.05);
      \draw[etrack] (-2.4,0.20) -- (0,0.20) -- (2.4,0.13);
      \draw[etrack] (-2.4,0.70) -- (0.07,0.70);
      \draw[eout] (0.07,0.70) -- (-0.97,1.30);
      \draw[etrack] (-2.4,-0.62) -- (0.10,-0.62);
      \draw[eout] (0.10,-0.62) -- (0.60,-1.49);
      \node[qk, fill=hscseriesC!85!black] at (0.30,0.70) {d};
      \node[qk, fill=hscseriesB] at (-0.45,-0.20) {u};
      \node[qk, fill=hscseriesB] at (0.35,-0.62) {u};
      \node[font=\footnotesize, text=hscforce!80!black] at (0,-1.75)
        {SLAC: far more scatter at large angles};
    \end{scope}
  \end{tikzpicture}
  \caption{Deep inelastic scattering: two models of the proton. The negative electrons are
    pulled slightly towards the positive charge. Model A gives only small deflections.
    SLAC saw large-angle bounces, which need point-like charges inside (Model B). Today we
    know these are quarks: u = up quark, d = down quark (\cref{ch5:sec:sm}).}\label{ch5:fig:dis}
\end{figure}

Why such fast electrons? A wave can only show details bigger than its wavelength. An
electron is also a wave (\cref{ch3:sec:debroglie}). To see inside a proton (diameter about
\qty{1.7e-15}{m}), its wavelength must be much smaller than that.

\begin{formulacard}{de Broglie wavelength}
  \cardformula{\lambda = \frac{h}{mv}}
  \begin{symbols}
    \sym{$\lambda$}{wavelength of the moving particle}{\si{m}}
    \sym{$h$}{Planck constant, \num{6.626e-34}}{\si{J.s}}
    \sym{$m$, $v$}{mass and speed of the particle}{\si{kg}, \si{m.s^{-1}}}
  \end{symbols}
  \inwords{more momentum ($mv$) means a shorter wavelength, so smaller details show up.}
  \whentouse{any moving particle with mass (electron, proton, neutron).}
  \whennottouse{photons ($E = hf$ instead). Above about $0.1c$, use the relativistic
    momentum $p_\mathrm{v}$ in place of $mv$.}
  \rearranged{$v = \dfrac{h}{m\lambda}$ \qquad $mv = \dfrac{h}{\lambda}$}
  \tensecond{electron at \qty{1.0e6}{m.s^{-1}}:
    $\lambda = \dfrac{\num{6.626e-34}}{\num{9.109e-31} \times \num{1.0e6}} = \qty{7.3e-10}{m}$.}
  \trap{use the electron's mass (\qty{9.109e-31}{kg}), not the proton's.}
\end{formulacard}

\begin{workedexample}{A slow electron cannot see inside a proton}\label{ch5:we:slow}
An electron moves at \qty{2.0e7}{m.s^{-1}}. Find its de Broglie wavelength. Compare it
with a proton's diameter, \qty{1.7e-15}{m}.
\given{$v = \qty{2.0e7}{m.s^{-1}}$ (only $0.07c$, well below $0.1c$, so
  $\lambda = \dfrac{h}{mv}$ works), \ $m = \qty{9.109e-31}{kg}$, \
  proton diameter $d = \qty{1.7e-15}{m}$}
\find{$\lambda$, then $\lambda \div d$}
\formula{$\lambda = \dfrac{h}{mv}$}
\substitute{\eqlines{\lambda &= \frac{\qty{6.626e-34}{J.s}}{\qty{9.109e-31}{kg} \times
  \qty{2.0e7}{m.s^{-1}}} = \qty{3.637e-11}{m}}\par
  \eqlines{\frac{\lambda}{d} &= \frac{\qty{3.637e-11}{m}}{\qty{1.7e-15}{m}}
  = \frac{3.637}{1.7} \times 10^{-11-(-15)} \quad\text{(divide the numbers)}\\
  &= 2.14 \times 10^{4} \quad\text{(subtract the powers: $-11 - (-15) = 4$)}}}
\answer{\final{$\lambda = \qty{3.6e-11}{m}$}: about \num{21000} times bigger than a
  proton, so it shows nothing inside one.}
\checkstep{$\lambda$ is about a third of the size of an atom (about \qty{e-10}{m}). Slow
  electrons like this are used to study atoms in crystals, which fits. (Powers of ten:
  see the toolkit in \cref{ch1:sec:bigbang}.)}
\end{workedexample}

To see inside a proton, the electron needs far more momentum. Near the speed of light,
momentum is no longer $mv$. You met the correct formula in Module 7.

\begin{formulacard}{Relativistic momentum (Module 7)}
  \cardformula{p_\mathrm{v} = \frac{m_0 v}{\sqrt{\left(1 - \dfrac{v^2}{c^2}\right)}}}
  \begin{symbols}
    \sym{$p_\mathrm{v}$}{momentum of a particle moving at speed $v$}{\si{kg.m.s^{-1}}}
    \sym{$m_0$}{rest mass (the mass measured when the particle is not moving)}{\si{kg}}
    \sym{$v$}{speed of the particle}{\si{m.s^{-1}}}
    \sym{$c$}{speed of light, \num{3.00e8}}{\si{m.s^{-1}}}
  \end{symbols}
  \inwords{near the speed of light, momentum grows much faster than speed.}
  \whentouse{a particle faster than about $0.1c$. Put $p_\mathrm{v}$ in place of $mv$:
    $\lambda = \dfrac{h}{p_\mathrm{v}}$ for its wavelength, $r = \dfrac{p_\mathrm{v}}{qB}$
    for its radius in a magnetic field. (These two forms are not printed on the sheet.)}
  \whennottouse{slow particles: below $0.1c$ the square root is 0.995 or more, so
    $p = mv$ is accurate to 0.5\%. Not for photons (they have no rest mass).}
  \rearranged{$m_0 = \dfrac{p_\mathrm{v}}{v}\sqrt{1 - \dfrac{v^2}{c^2}}$ \qquad
    for small $v$: $p_\mathrm{v} \approx m_0 v$}
  \tcbbreak % layout QA (book build): page break here, so "(continued)" is not a lone Trap line
  \tensecond{proton at $0.60c$: $m_0 v = \num{1.673e-27} \times 0.60 \times \num{3.00e8}
    = \num{3.01e-19}$; \ $\sqrt{1 - 0.60^2} = 0.80$; \ so
    $p_\mathrm{v} = \qty{3.8e-19}{kg.m.s^{-1}}$.}
  \trap{$m_0$ is the rest mass from the data sheet. Square $\dfrac{v}{c}$ \emph{before}
    you subtract it from 1.}
\end{formulacard}

\begin{workedexample}{Is 0.9999$c$ fast enough?}\label{ch5:we:fast}
An electron moves at $0.9999c$. Find its relativistic momentum and its wavelength. Can it
show structure inside a proton (diameter \qty{1.7e-15}{m})?
\given{$v = 0.9999c = 0.9999 \times \qty{3.00e8}{m.s^{-1}}$, \ rest mass
  $m_0 = \qty{9.109e-31}{kg}$, \ $d = \qty{1.7e-15}{m}$}
\find{$p_\mathrm{v}$, then $\lambda$, then $\lambda \div d$}
\formula{$p_\mathrm{v} = \dfrac{m_0 v}{\sqrt{\left(1 - \dfrac{v^2}{c^2}\right)}}$
  (on the sheet), \ then $\lambda = \dfrac{h}{p_\mathrm{v}}$ (not printed in this form:
  $\lambda = \dfrac{h}{mv}$ with $p_\mathrm{v}$ in place of $mv$)}
\substitute{\eqlines{m_0 v &= \qty{9.109e-31}{kg} \times 0.9999 \times \qty{3.00e8}{m.s^{-1}}\\
  &= \qty{2.732e-22}{kg.m.s^{-1}} \quad\text{(the top of the fraction)}\\
  \frac{v^2}{c^2} &= 0.9999^2 = 0.99980001 \quad\text{(now the bottom)}\\
  1 - \frac{v^2}{c^2} &= 1 - 0.99980001 = 0.00019999\\
  \sqrt{0.00019999} &= 0.014142 \quad\text{(square root last)}}\par
  \eqlines{p_\mathrm{v} &= \frac{\qty{2.732e-22}{kg.m.s^{-1}}}{0.014142} = \qty{1.932e-20}{kg.m.s^{-1}}\\
  \lambda &= \frac{\qty{6.626e-34}{J.s}}{\qty{1.932e-20}{kg.m.s^{-1}}} = \qty{3.429e-14}{m}\\
  \frac{\lambda}{d} &= \frac{\qty{3.429e-14}{m}}{\qty{1.7e-15}{m}} = 20}}
\answer{\final{$p_\mathrm{v} = \qty{1.93e-20}{kg.m.s^{-1}}$}, \
  \final{$\lambda = \qty{3.43e-14}{m}$}. No: $\lambda$ is still about 20 times
  \emph{bigger} than a proton.}
\checkstep{the speed is 15 times bigger than in \cref{ch5:we:slow}, but $p_\mathrm{v}$ is
  about 1060 times bigger (relativity). So $\lambda$ should be about 1060 times shorter:
  $\num{3.637e-11} \div 1060 = \qty{3.43e-14}{m}$ \cmark.}
\end{workedexample}

SLAC's \qty{20}{GeV} electrons had about 550 times more momentum than this. Their
wavelength was about \qty{6e-17}{m}, about 27 times \emph{smaller} than a proton. That
is why SLAC could see inside it.

\subsection{Evidence 2 and 3: beta decay and the particle zoo}

\begin{itemize}
  \item \textbf{Beta decay.} A free neutron decays:
        $\mathrm{n} \to \mathrm{p} + \mathrm{e}^- + \chFiveanu$ (\cref{ch4:sec:decay}). In the
        quark model, a neutron is udd: one up quark (u) and two down quarks (d). A proton
        is uud. So in this decay, just \emph{one} down quark changes into an up quark. The
        other two stay the same. This is easy to explain if the neutron has parts. So it
        supports the idea that protons and neutrons are not fundamental. (It is not proof
        on its own. The muon, a fundamental particle like a heavy electron, also decays.)
  \item \textbf{The particle zoo.} In the 1950s and 1960s, accelerators made hundreds of
        new particles. Scientists called this huge, confusing list the \term{particle zoo}.
        In 1964 Gell-Mann and Zweig proposed a simple idea. All these particles are built
        from just three kinds of quark: up, down and strange. Gell-Mann's earlier pattern
        of particles (1962) had also predicted a new particle, the omega-minus
        ($\Omega^-$). It was found in 1964, with the predicted charge and mass. In the
        quark model it is sss.
\end{itemize}

\subsection{Other subatomic particles}

Many other particles exist (\cref{ch5:tab:particles}). An \term{antiparticle} is a partner
with the same mass but the opposite charge. Its other built-in properties are opposite
too. So even a neutral particle can have a different antiparticle: an antineutron is not
a neutron. The \term{positron} ($\mathrm{e}^+$) is the antiparticle of the electron.

The positron was the first antiparticle found (\cref{ch5:fig:anderson}). Anderson saw it
in a \term{cloud chamber}: a box where fast charged particles leave a trail of tiny drops.
It came from \term{cosmic rays}: fast particles from space and the particles they make in
the air.

\begin{sidebyside}[0.40]
  \centering
  \begin{tikzpicture}[x=0.85cm, y=0.85cm]
    \draw[hscink, fill=hscmuted!8, line width=0.6pt] (-0.4,-0.3) rectangle (3.2,3.7);
    \foreach \p in {(0.1,0.5), (1.3,0.5), (0.1,1.4), (1.3,1.4), (0.1,2.7), (2.7,2.7),
                    (0.1,3.35), (2.7,3.35)}
      \node[text=hscfield, font=\scriptsize] at \p {$\times$};
    \fill[hscink!55] (-0.4,1.9) rectangle (3.2,2.15);
    \draw[hscvel, line width=0.9pt, midarrow] (2.6,0) arc[start angle=0, end angle=22.33,
      radius=5.0];
    \draw[hscvel, line width=0.9pt, -{Stealth[length=2.2mm, width=1.7mm]}] (2.122,2.15)
      arc[start angle=22.33, end angle=75, radius=2.0];
    \node[font=\scriptsize, anchor=west] at (3.3,2.02) {lead plate};
    \node[font=\scriptsize, anchor=west, align=left] at (3.3,0.8) {before:\\[-2pt]bends less};
    \node[font=\scriptsize, anchor=west, align=left] at (3.3,3.15) {after:\\[-2pt]bends more};
    \node[font=\scriptsize, text=hscfield, anchor=north] at (1.4,-0.4)
      {$\times$: $B$ into the page};
  \end{tikzpicture}
  \captionof{figure}{Sketch of Anderson's positron track (1932). It bends more after the
    lead plate, so it moved upwards.}\label{ch5:fig:anderson}
\nextside
  Anderson's cloud chamber (1932) sat inside a magnetic field.
  \begin{itemize}
    \item A lead plate slowed the particle, so its track curved more after the plate
          (less momentum, tighter curve). So it moved upwards.
    \item It moved up and bent left, with $B$ into the page. So its charge was
          \emph{positive}. An electron would bend right.
    \item The track was thin and long, like an electron's. A proton bent this much would
          stop much sooner. So its mass was about the electron's.
    \item No known particle was light and positive, so it was new.
  \end{itemize}
\end{sidebyside}

How did scientists know that each of these particles was new? In each case they saw a
mass, a charge or a behaviour that no known particle had (\cref{ch5:tab:particles}).

\begin{table}
  \caption{Some particles that are not protons, neutrons or electrons.}\label{ch5:tab:particles}
  \centering
  \small
  \begin{tabular}{@{}L{22mm}L{10mm}L{29mm}L{77mm}@{}}
    \toprule
    \colhead{Particle} & \colhead{Year} & \colhead{What it is} &
      \colhead{What was seen, and what it shows}\\
    \midrule
    positron $\mathrm{e}^+$ & 1932 & antiparticle of the electron & A cloud-chamber
      track (\cref{ch5:fig:anderson}). It had an electron's mass but bent the opposite
      way, so it was positive. Dirac's theory (1928) had predicted it; he named the
      anti-electron in 1931.\\
    muon $\mu^-$ & 1936 & a heavy copy of the electron & Cosmic-ray tracks. Charge $-e$,
      but about 207 times the electron's mass. No known particle fitted.\\
    pions $\pi^+$, $\pi^-$ & 1947 & particles made of quarks (mesons) & Cosmic-ray
      tracks in photographic film taken up mountains. Mass about 270 electron masses, between the
      muon and the proton. No known particle had this mass.\\
    antiproton $\overline{\mathrm{p}}$ & 1955 & antiparticle of the proton & Made in a
      proton accelerator (Bevatron, USA). It had the proton's mass but charge $-e$.\\
    neutrino $\nu$ & 1956 & a neutral particle with a tiny mass & In beta decay, electrons
      come out with a spread of energies (\cref{ch4:fig:betaspectrum}). Some energy seemed
      to vanish. Pauli (1930) said a hidden neutral particle carries it. It was detected
      near a nuclear reactor in 1956.\\
    \bottomrule
  \end{tabular}
\end{table}

\begin{mistakes}
\begin{itemize}
  \item \xmark\ ``A neutron is a proton plus an electron.'' \cmark\ The electron and
        antineutrino are \emph{created} in the decay.
  \item \xmark\ Using $\lambda = \dfrac{h}{m_0 v}$ near $c$. \cmark\ Use $p_\mathrm{v}$.
  \item \xmark\ ``The electron is made of quarks.'' \cmark\ It is a fundamental lepton.
  \item \xmark\ ``An antiparticle only has the opposite charge.'' \cmark\ Its other
        built-in properties are opposite too, so an antineutrino differs from a neutrino.
\end{itemize}
\end{mistakes}

\begin{examtip}
``Analyse the evidence that protons and neutrons are not fundamental'' (4--5 marks).
Link each observation to what it shows. SLAC electrons bounced at large angles far more
often than a smooth proton allows. So the proton has point-like charged parts (quarks).
Then add supporting evidence: beta decay ($\mathrm{d} \to \mathrm{u}$) and the quark
pattern of the particle zoo (the $\Omega^-$ prediction). For the other half of this dot
point, name each new particle. Say what was observed: for example, a positive track with an
electron's mass, or the spread of beta energies. Then say why that means a new particle
exists.
\end{examtip}

\begin{quickcheck}
  \item State what a fundamental particle is. Is the electron one?
  \item Electrons in a beam each have momentum \qty{8.0e-18}{kg.m.s^{-1}}. Find their
        wavelength. Can they show structure inside a proton (\qty{1.7e-15}{m})?
  \item Explain how beta-minus decay supports the idea that a neutron is not fundamental.
  \item Name two particles other than protons, neutrons and electrons, with the evidence
        for each.
\end{quickcheck}

%% ==================================================================
\section{The Standard Model: particles and forces}\label{ch5:sec:sm}
\covers{M8.IQ5.02}

\begin{bigidea}
The Standard Model is our best theory of what matter is made of. It has 12 matter
particles: 6 quarks and 6 leptons. Other particles, called bosons, carry the forces.
Quarks always join in groups to make hadrons, such as protons and neutrons.
\end{bigidea}

\begin{keywords}
  \kw{Quark}{a fundamental matter particle (fermion) that feels the strong force (the
    force that holds the nucleus together).}
  \kw{Lepton}{a fundamental matter particle that does not feel the strong force.}
  \kw{Hadron}{a particle made of quarks, such as a proton.}
  \kw{Fermion / boson}{a particle with half-number spin / with whole-number spin (spin is
    a built-in amount of turning).}
\end{keywords}

The \term{Standard Model} is the theory that lists the fundamental particles and the
forces between them (\cref{ch5:fig:sm}). It has passed nearly every test since the 1970s.

Every particle has \term{spin}: a built-in amount of turning, like a tiny spinning top. A
\term{fermion} has half-number spin ($\tfrac12$, $\tfrac32$, \ldots). All quarks and
leptons are fermions. A \term{boson} has whole-number spin (0, 1, 2, \ldots). The force
carriers and the Higgs are bosons.

\subsection{Quarks and leptons}

\begin{itemize}
  \item A \term{quark} is a fundamental matter particle (a fermion) that feels the strong
        force. There are six \term{flavours} (types) in three \term{generations}: each
        generation is a heavier copy of the one before. Ordinary matter uses only u and d.
  \item Up-type quarks (u, c, t) have charge $+\tfrac23 e$. Down-type quarks (d, s, b)
        have $-\tfrac13 e$. Quarks carry \term{colour charge}: the ``charge'' of the
        strong force (not real colour).
  \item A quark is never found alone. This is \term{confinement}: quarks are always locked
        inside hadrons.
  \item A \term{lepton} is a fundamental matter particle that does not feel the strong
        force. The electron, muon and tau have charge $-1e$. The three \term{neutrinos}
        have a tiny mass. They have no electric charge and no colour charge, so only the
        weak force and gravity act on them.
  \item Every particle has an antiparticle with the same mass and the opposite charge. (A
        few neutral particles, such as the photon, are their own antiparticle.) An
        antiquark has a bar: $\chFiveaq{u}$ has charge $-\tfrac23 e$ and $\chFiveaq{d}$
        has $+\tfrac13 e$.
\end{itemize}

\begin{figure}
  \centering
  \begin{tikzpicture}[x=1cm, y=1cm,
      cell/.style={draw=#1!75!black, fill=#1!14, rounded corners=1.5mm,
        minimum width=1.9cm, minimum height=1.36cm, line width=0.6pt,
        align=center, inner sep=1pt, execute at begin node={\setstretch{0.9}}},
      bcell/.style={cell=#1, minimum width=2.25cm}]
    \foreach \x/\g in {0/1, 2.05/2, 4.1/3}
      \node[font=\footnotesize\bfseries] at (\x,1.0) {Generation \g};
    \node[font=\footnotesize\bfseries] at (6.55,1.0) {Force carriers};
    \node[font=\footnotesize\bfseries] at (8.95,1.0) {Higgs};
    \node[cell=hscseriesB] at (0,0)    {{\large $\mathrm{u}$}\\{\scriptsize up}\\{\scriptsize $+\tfrac23$}};
    \node[cell=hscseriesB] at (2.05,0) {{\large $\mathrm{c}$}\\{\scriptsize charm}\\{\scriptsize $+\tfrac23$}};
    \node[cell=hscseriesB] at (4.1,0)  {{\large $\mathrm{t}$}\\{\scriptsize top}\\{\scriptsize $+\tfrac23$}};
    \node[cell=hscseriesB] at (0,-1.48)    {{\large $\mathrm{d}$}\\{\scriptsize down}\\{\scriptsize $-\tfrac13$}};
    \node[cell=hscseriesB] at (2.05,-1.48) {{\large $\mathrm{s}$}\\{\scriptsize strange}\\{\scriptsize $-\tfrac13$}};
    \node[cell=hscseriesB] at (4.1,-1.48)  {{\large $\mathrm{b}$}\\{\scriptsize bottom}\\{\scriptsize $-\tfrac13$}};
    \node[cell=hscseriesC] at (0,-2.96)    {{\large $\mathrm{e}^-$}\\{\scriptsize electron}\\{\scriptsize $-1$}};
    \node[cell=hscseriesC] at (2.05,-2.96) {{\large $\mu^-$}\\{\scriptsize muon}\\{\scriptsize $-1$}};
    \node[cell=hscseriesC] at (4.1,-2.96)  {{\large $\tau^-$}\\{\scriptsize tau}\\{\scriptsize $-1$}};
    \node[cell=hscseriesC] at (0,-4.44)    {{\large $\nu_\mathrm{e}$}\\{\scriptsize e-neutrino}\\{\scriptsize $0$}};
    \node[cell=hscseriesC] at (2.05,-4.44) {{\large $\nu_\mu$}\\{\scriptsize $\mu$-neutrino}\\{\scriptsize $0$}};
    \node[cell=hscseriesC] at (4.1,-4.44)  {{\large $\nu_\tau$}\\{\scriptsize $\tau$-neutrino}\\{\scriptsize $0$}};
    \node[bcell=hscseriesA] at (6.55,0)     {{\large $\mathrm{g}$}\\{\scriptsize gluon}\\{\scriptsize strong force}};
    \node[bcell=hscseriesA] at (6.55,-1.48) {{\large $\gamma$}\\{\scriptsize photon}\\{\scriptsize EM force}};
    \node[bcell=hscseriesA] at (6.55,-2.96) {{\large $\mathrm{W}^\pm$}\\{\scriptsize W bosons}\\{\scriptsize weak force}};
    \node[bcell=hscseriesA] at (6.55,-4.44) {{\large $\mathrm{Z}^0$}\\{\scriptsize Z boson}\\{\scriptsize weak force}};
    \node[cell=hscbrand] at (8.95,0) {{\large $\mathrm{H}$}\\{\scriptsize Higgs boson}\\{\scriptsize spin 0}};
    \node[draw=hscmuted, dashed, rounded corners=1.5mm, minimum width=1.9cm,
      minimum height=1.36cm, align=center, inner sep=1pt, line width=0.6pt,
      execute at begin node={\setstretch{0.9}}] (grav) at (8.95,-2.35)
      {{\large\color{hscmuted} G?}\\{\scriptsize\color{hscmuted} graviton}};
    \node[font=\scriptsize, text=hscmuted, align=center, below=1pt of grav]
      {never found;\\ NOT in the\\ Standard Model};
    \draw[hscseriesB!75!black, line width=0.8pt] (-1.15,0.68) -- (-1.25,0.68) -- (-1.25,-2.16) -- (-1.15,-2.16);
    \node[rotate=90, font=\footnotesize\bfseries, text=hscseriesB!75!black] at (-1.55,-0.74) {QUARKS};
    \draw[hscseriesC!75!black, line width=0.8pt] (-1.15,-2.28) -- (-1.25,-2.28) -- (-1.25,-5.12) -- (-1.15,-5.12);
    \node[rotate=90, font=\footnotesize\bfseries, text=hscseriesC!75!black] at (-1.55,-3.70) {LEPTONS};
    \draw[hscmuted, line width=0.6pt] (-0.95,-5.3) -- (5.05,-5.3);
    \node[font=\scriptsize] at (2.05,-5.57) {matter particles: fermions (spin $\tfrac12$)};
    \draw[hscmuted, line width=0.6pt] (5.42,-5.3) -- (10.1,-5.3);
    \node[font=\scriptsize] at (7.75,-5.57) {bosons (whole-number spin)};
  \end{tikzpicture}
  \caption{The particles of the Standard Model. Numbers are charges in units of $e$.
    EM = electromagnetic.}\label{ch5:fig:sm}
\end{figure}

\subsection{The four fundamental forces}

Imagine two people on skateboards passing a heavy medicine ball. The thrower rolls
backwards. The catcher rolls backwards too. Each pass pushes them further apart. In the
same way, two electrons swap photons, and this pushes them apart. (This picture only
explains pushing apart. Attraction needs the full theory.) A boson that carries a force
like this is an \term{exchange particle} (\cref{ch5:tab:forces}).

\begin{table}
  \caption{The four forces. Strengths are rough values compared with the strong force
    (books differ a little).}\label{ch5:tab:forces}
  \centering
  \small
  \begin{tabular}{@{}L{17mm}L{22mm}L{26mm}L{13mm}L{14mm}L{36mm}@{}}
    \toprule
    \colhead{Force} & \colhead{Carrier} & \colhead{Acts on} & \colhead{Strength}
      & \colhead{Range} & \colhead{What it does}\\
    \midrule
    strong & gluons (g) & quarks and gluons & 1 & \qty{e-15}{m} & holds quarks in hadrons
      and holds the nucleus together\\
    electro\-magnetic & photon ($\gamma$) & charged particles & $10^{-2}$ & no limit
      & holds atoms together; light\\
    weak & $\mathrm{W}^+$, $\mathrm{W}^-$, $\mathrm{Z}^0$ & all quarks and leptons
      & $10^{-6}$ & \qty{e-18}{m} & beta decay; changes quark flavour\\
    gravity & graviton? (never found) & everything with mass or energy & $10^{-38}$
      & no limit & planets, stars, galaxies\\
    \bottomrule
  \end{tabular}
\end{table}

\begin{itemize}
  \item \textbf{Strong force.} The \term{strong force} (carried by gluons) holds quarks
        together inside protons and neutrons. A little of this pull leaks out of each
        proton and neutron. This left-over pull holds the nucleus together, even though
        its protons repel each other.
  \item \textbf{Weak force.} The \term{weak force} is the only force that can change a
        quark's flavour (type). It is carried by three bosons: $\mathrm{W}^+$,
        $\mathrm{W}^-$ and $\mathrm{Z}^0$. They are very heavy: about 86 and 97 times the
        mass of a proton. A heavy exchange particle can only exist for a very short time.
        So it cannot travel far, and the weak force has a tiny range (about
        \qty{e-18}{m}). \Cref{ch5:fig:beta} shows how it changes a quark in beta decay.
\end{itemize}

\begin{sidebyside}[0.42]
  \centering
  \begin{tikzpicture}[x=0.95cm, y=0.95cm,
      had/.style={draw=hscbrand!60, fill=hscbrand!6, line width=0.7pt},
      qk/.style={circle, text=white, font=\scriptsize\bfseries, minimum size=0.46cm,
        inner sep=0pt},
      hl/.style={draw=hscforce, line width=1.3pt}]
    \draw[had] (0,0) circle (0.8);
    \node[qk, fill=hscseriesB] at (-0.3,0.25) {u};
    \node[qk, fill=hscseriesC!85!black] at (0.3,0.25) {d};
    \node[qk, fill=hscseriesC!85!black, hl] at (0,-0.3) {d};
    \node[font=\scriptsize\bfseries] at (0,1.1) {neutron (udd)};
    \draw[hscink, line width=0.9pt, -{Stealth[length=2.2mm, width=1.7mm]}] (0.95,0) -- (2.45,0);
    \node[font=\scriptsize, text=hscforce, anchor=south] at (1.7,0.05) {d $\to$ u};
    \draw[had] (3.4,0) circle (0.8);
    \node[qk, fill=hscseriesB] at (3.1,0.25) {u};
    \node[qk, fill=hscseriesC!85!black] at (3.7,0.25) {d};
    \node[qk, fill=hscseriesB, hl] at (3.4,-0.3) {u};
    \node[font=\scriptsize\bfseries] at (3.4,1.1) {proton (uud)};
    \draw[hscacc, line width=1pt, decorate, decoration={snake, amplitude=0.8mm,
      segment length=2.2mm}] (1.7,-0.05) -- (1.7,-1.1);
    \node[font=\scriptsize, text=hscacc!80!black, anchor=west] at (1.8,-0.6) {$\mathrm{W}^-$};
    \draw[hscvel, line width=0.9pt, -{Stealth[length=2.2mm, width=1.7mm]}] (1.7,-1.1) -- (0.75,-1.85);
    \draw[hscvel, line width=0.9pt, -{Stealth[length=2.2mm, width=1.7mm]}] (1.7,-1.1) -- (2.65,-1.85);
    \node[font=\scriptsize, anchor=north, align=center] at (0.6,-1.9) {electron\\[-2pt]$\mathrm{e}^-$};
    \node[font=\scriptsize, anchor=north, align=center] at (2.85,-1.9) {antineutrino\\[-2pt]$\chFiveanu$};
  \end{tikzpicture}
  \captionof{figure}{Beta-minus decay at the quark level. The ringed quark is the one
    that changes.}\label{ch5:fig:beta}
\nextside
  \begin{itemize}
    \item \textbf{Beta-minus decay:} a down quark becomes an up quark and gives out a
          $\mathrm{W}^-$: $\mathrm{d} \to \mathrm{u} + \mathrm{W}^-$. The $\mathrm{W}^-$
          quickly becomes an electron and an antineutrino:
          $\mathrm{W}^- \to \mathrm{e}^- + \chFiveanu$. So the neutron (udd) becomes a
          proton (uud).
    \item \textbf{Beta-plus decay} (only inside a nucleus):
          $\mathrm{u} \to \mathrm{d} + \mathrm{W}^+$, then
          $\mathrm{W}^+ \to \mathrm{e}^+ + \nu_\mathrm{e}$. A free proton cannot do this:
          it is lighter than a neutron.
    \item Charge balances at every step: $-\tfrac13 = +\tfrac23 + (-1)$.
  \end{itemize}
\end{sidebyside}

\begin{itemize}
  \item \textbf{The Higgs field} fills all space. Most fundamental particles get their mass
        by interacting with it. (How neutrinos get their tiny mass is not yet known.) The
        photon and gluon do not interact with it, so they have no mass. A rough picture:
        it is harder to get moving in water than in air. In the same way, a particle that
        interacts strongly with the Higgs field is harder to speed up, so it has more
        mass. (This is only a picture: the field does not slow particles down like
        friction.) A ripple in the field is the \term{Higgs boson}, found in 2012. Most
        of a proton's mass is not from the Higgs: about 99\% is the energy of its moving
        quarks and gluons.
\end{itemize}

\subsection{Hadrons: baryons and mesons}

\begin{sidebyside}[0.40]
  \centering
  \begin{tikzpicture}[x=0.85cm, y=0.85cm,
      had/.style={draw=hscbrand!60, fill=hscbrand!6, line width=0.7pt},
      upq/.style={circle, fill=hscseriesB!80, text=white, font=\small\bfseries,
        minimum size=0.48cm, inner sep=0pt},
      dnq/.style={circle, fill=hscseriesC!85!black, text=white, font=\small\bfseries,
        minimum size=0.48cm, inner sep=0pt},
      glue/.style={decorate, decoration={coil, aspect=0, segment length=1.6mm,
        amplitude=0.6mm}, hscmuted, line width=0.5pt}]
    \draw[had] (0,0) circle (0.95);
    \node[upq] (a) at (-0.38,0.30) {u}; \node[upq] (b) at (0.40,0.30) {u};
    \node[dnq] (c) at (0,-0.38) {d};
    \draw[glue] (a) -- (b); \draw[glue] (b) -- (c); \draw[glue] (c) -- (a);
    \node[font=\footnotesize\bfseries] at (0,1.25) {proton (baryon)};
    \node[font=\footnotesize] at (0,-1.3) {$+\tfrac23 + \tfrac23 - \tfrac13 = +1$};
    \begin{scope}[xshift=2.9cm]
      \draw[had] (0,0) ellipse (0.95 and 0.75);
      \node[upq] (a) at (-0.40,0) {u};
      \node[circle, draw=hscseriesC!85!black, line width=1pt, fill=white,
        text=hscseriesC!70!black, font=\small\bfseries, minimum size=0.48cm,
        inner sep=0pt] (b) at (0.40,0) {$\overline{\mathrm{d}}$};
      \draw[glue] (a) -- (b);
      \node[font=\footnotesize\bfseries] at (0,1.25) {pion $\pi^+$ (meson)};
      \node[font=\footnotesize] at (0,-1.3) {$+\tfrac23 + \tfrac13 = +1$};
    \end{scope}
  \end{tikzpicture}
  \captionof{figure}{A baryon and a meson. The springs stand for gluons.}\label{ch5:fig:hadrons}
\nextside
  A \term{hadron} is a particle made of quarks held together by the strong force
  (\cref{ch5:fig:hadrons}).
  \begin{itemize}
    \item A \term{baryon} is a hadron of three quarks: proton (uud), neutron (udd). Three
          half-number spins add to a half-number, so baryons are fermions.
    \item A \term{meson} is a hadron of a quark and an antiquark: pion $\pi^+$
          (u$\chFiveaq{d}$). Two half-number spins add to a whole number, so mesons are
          bosons. All mesons are unstable.
  \end{itemize}
\end{sidebyside}

\begin{toolkit}{Adding thirds}
Write every charge as a fraction over 3. Then add the top numbers:
$\tfrac23 + \tfrac23 - \tfrac13 = \dfrac{2 + 2 - 1}{3} = \dfrac{3}{3} = 1$. Keep each minus
sign with its top number. To get coulombs, multiply by $e = \qty{1.602e-19}{C}$.
\end{toolkit}

\begin{formulacard}[notonsheet]{Charge of a hadron}
  \cardformula{Q = q_1 + q_2 + q_3 \ \ \text{(baryon)} \qquad Q = q_1 + q_2 \ \ \text{(meson)}}
  \begin{symbols}
    \sym{$Q$}{total charge of the hadron}{\si{C} (or $e$)}
    \sym{$q_1, q_2, q_3$}{charge of each quark: u, c, t $+\tfrac23 e$; d, s, b
      $-\tfrac13 e$; antiquarks opposite}{\si{C} (or $e$)}
  \end{symbols}
  \inwords{add up the charges of the quarks inside.}
  \whentouse{any baryon or meson; or backwards, to find a missing quark.}
  \whennottouse{leptons: they are not made of quarks.}
  \rearranged{missing quark $= Q - {}$(sum of the known quarks); \ in coulombs, multiply
    the charge in $e$ by \qty{1.602e-19}{C}.}
  \tensecond{proton uud: $+\tfrac23 + \tfrac23 - \tfrac13 = +1e$.}
  \derive{charge is conserved and adds up: a group's charge is the sum of its parts'
    charges. The quark charges must be learnt.}
  \trap{flip the sign for an antiquark: $\chFiveaq{d}$ is $+\tfrac13 e$.}
\end{formulacard}

\begin{workedexample}{Charges of two hadrons}\label{ch5:we:charges}
The $\Xi^0$ (``xi-zero'') baryon is uss. The $\mathrm{K}^-$ (``K-minus'', a kaon) meson is
s$\chFiveaq{u}$. Find the charge of each, in units of $e$ and in coulombs.
\given{u $= +\tfrac23 e$, \ s $= -\tfrac13 e$, \ $\chFiveaq{u} = -\tfrac23 e$ (sign flipped)}
\find{$Q$ for $\Xi^0$ and for $\mathrm{K}^-$}
\formula{$Q = q_1 + q_2 + q_3$ (baryon), \ $Q = q_1 + q_2$ (meson)}
\substitute{\eqlines{Q(\Xi^0) &= +\tfrac23 - \tfrac13 - \tfrac13 = \dfrac{2 - 1 - 1}{3} = 0\\
  Q(\Xi^0) &= 0 \times \qty{1.602e-19}{C} = \qty{0}{C} \quad\text{(change to coulombs)}\\
  Q(\mathrm{K}^-) &= -\tfrac13 - \tfrac23 = \dfrac{-1 - 2}{3} = -1\\
  Q(\mathrm{K}^-) &= -1 \times \qty{1.602e-19}{C} = \qty{-1.602e-19}{C}}}
\answer{\final{$Q(\Xi^0) = 0 = \qty{0}{C}$} \ and \
  \final{$Q(\mathrm{K}^-) = -1e = \qty{-1.602e-19}{C}$}}
\checkstep{the small 0 and $-$ in the names $\Xi^0$ and $\mathrm{K}^-$ show their charges
  \cmark. A hadron's charge is always a whole number of $e$.}
\end{workedexample}

\begin{workedexample}{Counting quarks in a nucleus}\label{ch5:we:nitrogen}
A nitrogen-14 nucleus has 7 protons and 7 neutrons. How many up and down quarks does it
contain? Use them to check its charge.
\given{7 protons (uud) and $14 - 7 = 7$ neutrons (udd)}
\find{the number of u and d quarks, and the total charge}
\formula{count the quarks in each nucleon (proton or neutron); then add all the quark
  charges: $Q = \sum q$ ($\sum$ means ``add them all up'')}
\substitute{\eqlines{\text{up} &= 7 \times 2 + 7 \times 1 = 21 \quad\text{(2 per proton, 1 per neutron)}\\
  \text{down} &= 7 \times 1 + 7 \times 2 = 21\\
  Q &= 21 \times \left(+\tfrac23\right) + 21 \times \left(-\tfrac13\right)
  = \tfrac{42}{3} - \tfrac{21}{3} = 14 - 7 = +7}}
\answer{\final{21 up and 21 down quarks}, total charge $+7e$.}
\checkstep{7 protons give $+7e$ \cmark. $3 \times 14 = 42$ quarks \cmark. Electrons are
  leptons, so they add no quarks.}
\end{workedexample}

\begin{workedexample}{Finding a missing quark}\label{ch5:we:missing}
A baryon has charge $-1e$. Two of its quarks are s and d. Its third quark is from the
first generation (u or d). (a) Find the third quark. (b) Is this baryon a fermion or a
boson? Explain. (c) It decays when its s quark changes into a u quark. Which force does
this, and which boson is given out?
\given{$Q = -1e$; \ s $= -\tfrac13 e$, \ d $= -\tfrac13 e$; \ third quark u
  ($+\tfrac23 e$) or d ($-\tfrac13 e$)}
\find{(a) the third quark \ (b) fermion or boson \ (c) the force and the boson}
\formula{missing quark $= Q - {}$(sum of the known quarks); \ charge is conserved in
  every change}
\substitute{\eqlines{q_3 &= -1 - \left(-\tfrac13\right) - \left(-\tfrac13\right)
  \quad\text{(take away the two known quarks)}\\
  &= -\tfrac33 + \tfrac13 + \tfrac13 = -\tfrac13 \quad\text{(so it is d, not u)}\\
  \mathrm{s} &\to \mathrm{u} + \mathrm{W}^-: \quad -\tfrac13 = +\tfrac23 + (-1)
  \quad\text{(charge balances)}}}
\answer{(a) \final{d}: the baryon is dds (the sigma-minus, $\Sigma^-$). (b) Each quark has
  spin $\tfrac12$. Three half-number spins add to a half-number, so it is a
  \final{fermion}. (c) Only the \final{weak force} changes a quark's flavour. The s quark
  gives out a \final{$\mathrm{W}^-$}.}
\checkstep{dds: $-\tfrac13 - \tfrac13 - \tfrac13 = -1$ \cmark. With u instead,
  $+\tfrac23 - \tfrac13 - \tfrac13 = 0$, not $-1$, so u is ruled out \cmark.}
\end{workedexample}

\begin{mistakes}
\begin{itemize}
  \item \xmark\ Not flipping an antiquark's sign. \cmark\ $\chFiveaq{u} = -\tfrac23 e$,
        $\chFiveaq{d} = +\tfrac13 e$.
  \item \xmark\ ``Mesons are fermions.'' \cmark\ Quark + antiquark gives whole-number
        spin: mesons are bosons. Hadrons include both baryons and mesons.
  \item \xmark\ ``The graviton is part of the Standard Model.'' \cmark\ It has never been
        found, and gravity is not in the Standard Model.
  \item \xmark\ $\mathrm{n} \to \mathrm{p} + \mathrm{e}^- + \nu$. \cmark\ Beta-minus decay
        gives an \emph{anti}neutrino, $\chFiveanu$.
\end{itemize}
\end{mistakes}

\begin{examtip}
``Compare quarks and leptons'' (4 marks) needs two similarities and two differences,
each with an example. Same: both are fundamental fermions in 3 generations of 2, and both
feel the weak force. Different: only quarks feel the strong force; quark charges are
$+\tfrac23 e$ or $-\tfrac13 e$, lepton charges $-1e$ or 0. For beta decay, show the quark
change ($\mathrm{d} \to \mathrm{u}$), the $\mathrm{W}^-$ and the $\chFiveanu$.
\end{examtip}

\begin{quickcheck}
  \item The antiproton is $\chFiveaq{u}\,\chFiveaq{u}\,\chFiveaq{d}$. Find its charge in
        $e$ and in coulombs.
  \item Sort into lepton, baryon, meson or force carrier: muon, $\pi^-$, proton, photon,
        electron neutrino. Which are fundamental?
  \item A baryon has charge $+2e$ and only first-generation quarks. What is it made of?
  \item List the four forces from weakest to strongest, with the carrier of each.
  \item Write beta-plus decay at the quark level. Which boson is involved?
\end{quickcheck}

%% ==================================================================
\section{Particle accelerators: how they work and what they show}\label{ch5:sec:accel}
\covers{M8.IQ5.03}

\begin{bigidea}
Accelerators use electric fields to speed up charged particles and magnetic fields to
steer them. Fast particles show tiny details, and their collisions make new heavy
particles. This is how the Standard Model has been tested.
\end{bigidea}

\begin{keywords}
  \kw{Linac}{a linear accelerator: particles are pushed along a straight line.}
  \kw{Cyclotron}{particles spiral outwards in a fixed magnetic field.}
  \kw{Synchrotron}{a ring where the magnetic field grows as the particles gain momentum,
    so they stay on one circle.}
  \kw{Collider}{an accelerator where two beams hit each other head-on.}
\end{keywords}

A \term{particle accelerator} is a machine that makes charged particles move very fast.
There are two reasons to build one. First, more momentum gives a shorter wavelength
($\lambda = \dfrac{h}{mv}$), so smaller details show up. Second, in a collision, kinetic
energy can turn into the mass of new particles ($E = mc^2$). To make a particle and its
antiparticle, you need at least $2mc^2$ (\cref{ch1:we:proton}).

\begin{refresher}{Charges in fields (Module 6)}
An electric field pushes a positive charge along the field. It pushes a negative charge
(such as an electron) the opposite way. A magnetic force, $F = qv_\perp B = qvB\sin\theta$
($\theta$ is the angle between $v$ and $B$), is always at \ang{90} to the velocity. It changes the direction, not the speed, so it does
no work. Only \emph{charged} particles can be accelerated.
\end{refresher}

\begin{toolkit}{Electronvolts}
One electronvolt (\si{eV}, \cref{ch1:sec:bigbang}) is the energy a charge $e$ gains through \qty{1}{V}:
$\qty{1}{eV} = \qty{1.602e-19}{J}$. Prefixes: k $= 10^3$, M $= 10^6$, G $= 10^9$,
T $= 10^{12}$. A proton through \qty{250}{kV} gains \qty{250}{keV}; an alpha particle
($q = 2e$) gains twice as much. To change eV to J, multiply by \num{1.602e-19}.
\end{toolkit}

\begin{formulacard}{Energy gained by a charge}
  \cardformula{W = qV}
  \begin{symbols}
    \sym{$W$}{work done on the charge (= the kinetic energy it gains)}{\si{J}}
    \sym{$q$}{charge of the particle}{\si{C}}
    \sym{$V$}{potential difference (voltage) it moves through}{\si{V}}
  \end{symbols}
  \inwords{a charge that moves through a voltage $V$ gains energy $qV$.}
  \whentouse{a charge speeding up between plates, in a linac gap or a cyclotron gap. From
    rest, after $n$ equal gaps: $K = nqV$.}
  \whennottouse{magnetic fields (they do no work); neutral particles ($q = 0$). Do not use
    $\tfrac12 mv^2 = qV$ to find a speed above about $0.1c$.}
  \rearranged{$V = \dfrac{W}{q}$ \qquad from rest: $\tfrac12 mv^2 = qV$, so
    $v = \sqrt{\dfrac{2qV}{m}}$ (below about $0.1c$)}
  \tensecond{electron through \qty{5.0}{kV}:
    $W = \num{1.602e-19} \times \num{5.0e3} = \qty{8.0e-16}{J} = \qty{5.0}{keV}$.}
  \trap{change \si{kV} to \si{V} first. An alpha particle has $q = 2e$, not $e$, so it
    gains twice as many eV.}
\end{formulacard}

\subsection{The linear accelerator (linac)}

A \term{linac} is a row of hollow metal \term{drift tubes}: tubes the particles coast
through (move at a steady speed, with no push). Every second tube is joined to one end of
an AC supply. The tubes in between are joined to the other end (\cref{ch5:fig:linac}).

\begin{figure}
  \centering
  \begin{tikzpicture}[x=1cm, y=1cm,
      tube/.style={draw=hscink, fill=hscmuted!18, line width=0.7pt}]
    \draw[draw=hscink, fill=hscbrand!10, line width=0.7pt] (-1.75,-0.35) rectangle (-0.55,0.35);
    \node[font=\scriptsize, align=center] at (-1.15,0) {proton\\[-2pt]source};
    \foreach \a/\b/\s in {0/0.850/+, 1.200/2.402/-, 2.752/4.224/+, 4.574/6.274/-, 6.624/8.525/+}{
      \draw[tube] (\a,-0.3) rectangle (\b,0.3);
      \node[font=\small\bfseries] at ({(\a+\b)/2},0) {$\s$};
    }
    \draw[guide] (-0.55,0) -- (0,0);
    \draw[velocity] (8.6,0) -- (9.6,0) node[midway, below, font=\scriptsize, text=hscvel] {beam};
    \fill[hscvel] (1.02,0) circle (2.2pt);
    \draw[force] (1.0,0.55) -- (1.6,0.55);
    \node[font=\scriptsize, text=hscforce, anchor=south] at (1.3,0.62) {push};
    % odd tubes -> top wire; even tubes -> bottom wire. The AC supply sits above the
    % tubes, off the beam line; its second wire runs round the left end.
    \foreach \c in {0.425, 3.488, 7.575} \draw[hscink, line width=0.6pt] (\c,0.3) -- (\c,1.15);
    \foreach \c in {1.801, 5.424} \draw[hscink, line width=0.6pt] (\c,-0.3) -- (\c,-1.15);
    \draw[hscink, line width=0.6pt] (0.425,1.15) -- (7.575,1.15);
    \draw[hscink, line width=0.6pt] (4.6,1.15) -- (4.6,1.45);
    \draw[hscink, fill=white, line width=0.7pt] (4.6,1.75) circle (0.3);
    \node[font=\small] at (4.6,1.75) {$\sim$};
    \node[font=\scriptsize, anchor=west] at (4.95,1.75) {AC supply};
    \draw[hscink, line width=0.6pt] (4.6,2.05) -- (4.6,2.3) -- (-2.1,2.3) -- (-2.1,-1.15)
      -- (5.424,-1.15);
    \node[font=\scriptsize, anchor=north] at (3.488,-0.38) {no field inside a tube};
    \draw[hscmuted, line width=0.5pt] (6.45,0.62) -- (6.45,0.36);
    \node[font=\scriptsize, anchor=south] at (6.45,0.6) {gap: field here};
    \draw[displacement] (2.75,-1.55) -- (8.5,-1.55) node[midway, below, font=\scriptsize, text=hscmuted]
      {each tube is longer than the one before};
  \end{tikzpicture}
  \caption{A linac at one instant. Half a cycle later every $+$ and $-$ swaps, so the
    next gap also pushes the proton forward.}\label{ch5:fig:linac}
\end{figure}

There is no field inside a tube, so the particle coasts. In each gap the field pushes it
forward and it gains $qV$. While it is inside a tube, the AC swaps $+$ and $-$, so the
next gap is ready to push. The particle speeds up, so each tube is \emph{longer}. This
keeps the time inside every tube equal to half an AC cycle. Hospitals use electron linacs
to make X-rays that treat cancer. The SLAC linac is \qty{3.2}{km} long.

\begin{workedexample}{Designing a proton linac}\label{ch5:we:linac}
The first stage of a proton accelerator must give each proton \qty{2.00}{MeV}. The
protons start from rest. Each gap has \qty{250}{kV} across it. How many gaps are needed?
How fast is a proton at the end?
\given{$K = \qty{2.00}{MeV}$, \ $V = \qty{250}{kV} = 250 \times 10^3 = \qty{2.50e5}{V}$, \
  $q = e = \qty{1.602e-19}{C}$, \ $m = \qty{1.673e-27}{kg}$}
\find{the number of gaps $n$ and the final speed $v$}
\formula{$W = qV$ for each gap; from rest $K = nqV$; then $K = \tfrac12 mv^2$, so
  $v = \sqrt{\dfrac{2K}{m}}$}
\substitute{\eqlines{W &= e \times \qty{250}{kV} = \qty{250}{keV} = \qty{0.250}{MeV}
  \quad\text{(energy from one gap)}\\
  n &= \frac{\qty{2.00}{MeV}}{\qty{0.250}{MeV}} = 8 \quad\text{(number of gaps)}}\par
  \eqlines{K &= \num{2.00e6} \times \qty{1.602e-19}{J} = \qty{3.204e-13}{J}
  \quad\text{(eV to J)}\\
  v &= \sqrt{\frac{2 \times \qty{3.204e-13}{J}}{\qty{1.673e-27}{kg}}}
  = \sqrt{\num{3.830e14}} = \qty{1.957e7}{m.s^{-1}}}}
\answer{\final{$n = 8$ gaps}, \ \final{$v = \qty{1.96e7}{m.s^{-1}}$}}
\checkstep{$\dfrac{v}{c} = \dfrac{\num{1.957e7}}{\num{3.00e8}} = 0.065$, well below
  $0.1c$, so $\tfrac12 mv^2$ is fine (relativity changes $v$ by only 0.2\%).}
\end{workedexample}

\subsection{The cyclotron}

A \term{cyclotron} has two hollow, D-shaped electrodes called \term{dees}, with a gap
between them. A large magnet makes a uniform field (the same strength and direction
everywhere) through the dees. An AC voltage is connected across the gap
(\cref{ch5:fig:cyclotron}). The particles are \term{ions}: charged atoms (a proton is a
hydrogen ion).

\begin{sidebyside}[0.47]
  \centering
  \begin{tikzpicture}[x=1.3cm, y=1.3cm]
    \fill[hscseriesB!12] (0.15,-1.45) arc[start angle=-90, end angle=90, radius=1.75] -- cycle;
    \fill[hscseriesB!12] (-0.15,2.05) arc[start angle=90, end angle=270, radius=1.75] -- cycle;
    \draw[hscseriesB!75!black, line width=0.8pt] (0.15,-1.45) arc[start angle=-90, end angle=90, radius=1.75] -- cycle;
    \draw[hscseriesB!75!black, line width=0.8pt] (-0.15,2.05) arc[start angle=90, end angle=270, radius=1.75] -- cycle;
    \foreach \p in {(-1.45,1.25), (-1.45,-0.55), (-0.95,1.62), (-0.75,-1.25), (1.45,1.25),
                    (1.45,-0.65), (0.95,1.62), (0.75,-1.25), (-1.72,0.3)}
      \node[text=hscfield, font=\small] at \p {$\times$};
    \draw[hscvel, line width=0.9pt, postaction={decorate}, decoration={markings,
        mark=at position 0.12 with {\arrow{Stealth[length=2mm, width=1.6mm]}},
        mark=at position 0.45 with {\arrow{Stealth[length=2mm, width=1.6mm]}},
        mark=at position 0.80 with {\arrow{Stealth[length=2mm, width=1.6mm]}}}]
      (0,0) -- (0.15,0) arc[start angle=-90, end angle=90, radius=0.400] -- ++(-0.30,0)
      arc[start angle=90, end angle=270, radius=0.566] -- ++(0.30,0)
      arc[start angle=-90, end angle=90, radius=0.693] -- ++(-0.30,0)
      arc[start angle=90, end angle=270, radius=0.800] -- ++(0.30,0)
      arc[start angle=-90, end angle=90, radius=0.894] -- ++(-0.30,0)
      arc[start angle=90, end angle=270, radius=0.980] -- ++(0.30,0)
      arc[start angle=-90, end angle=90, radius=1.058] -- ++(-0.30,0)
      arc[start angle=90, end angle=270, radius=1.131] -- ++(0.30,0)
      arc[start angle=-90, end angle=90, radius=1.200] -- ++(-0.30,0)
      arc[start angle=90, end angle=270, radius=1.265] -- ++(0.30,0)
      arc[start angle=-90, end angle=-15, radius=1.327];
    \draw[velocity] (1.431,-0.009) -- (1.846,1.536);
    \node[font=\scriptsize, text=hscvel, anchor=west] at (1.86,1.5) {exit};
    \fill[hscink] (0,0) circle (1.6pt);
    \draw[hscink, line width=0.6pt] (-0.6,1.98) -- (-0.6,2.35) -- (-0.25,2.35);
    \draw[hscink, line width=0.6pt] (0.6,1.98) -- (0.6,2.35) -- (0.25,2.35);
    \draw[hscink, fill=white, line width=0.7pt] (0,2.35) circle (0.25);
    \node[font=\small] at (0,2.35) {$\sim$};
    \node[font=\scriptsize, anchor=west] at (0.62,2.42) {AC};
    \node[font=\scriptsize, text=hscseriesB!65!black] at (-1.15,-1.62) {dee 1};
    \node[font=\scriptsize, text=hscseriesB!65!black] at (1.15,-1.62) {dee 2};
    \node[font=\scriptsize] at (0,-1.62) {gap};
    \node[font=\scriptsize, align=center, anchor=north] at (0,-1.8)
      {gap: the AC pushes the ion here\\[-1pt]$\times$: $B$ into the page \quad
       dot: ion source};
  \end{tikzpicture}
  \captionof{figure}{Top view of a cyclotron. A positive ion starts at the centre and
    turns anticlockwise.}\label{ch5:fig:cyclotron}
\nextside
  \begin{itemize}
    \item Each time the ion crosses the gap, the electric field speeds it up.
    \item Inside a dee there is no electric field. The magnetic force bends the ion
          around a half circle.
    \item Faster means a bigger circle, so the path spirals outwards.
    \item The AC flips while the ion is inside a dee, so every crossing is a push.
    \item At the edge the fast ions are steered out.
  \end{itemize}
\end{sidebyside}

\begin{formulacard}[notonsheet]{A charge moving in a circle in a magnetic field}
  \cardformula{r = \frac{mv}{qB} \qquad\qquad T = \frac{2\pi m}{qB}}
  \begin{symbols}
    \sym{$r$, $T$}{radius of the circle; period (time for one full circle)}{\si{m}, \si{s}}
    \sym{$m$, $v$}{mass and speed (speed at \ang{90} to $B$)}{\si{kg}, \si{m.s^{-1}}}
    \sym{$q$, $B$}{charge; magnetic field strength}{\si{C}, \si{T}}
  \end{symbols}
  \inwords{faster or heavier particles curve less, but the time for one circle does not
    depend on speed: a faster particle goes round a bigger circle in the same time.}
  \whentouse{a charge moving at \ang{90} to a uniform field at speeds well below $c$:
    cyclotron, slow tracks in detectors. A cyclotron's AC has one cycle per turn, so
    $f = \dfrac{1}{T}$.}
  \whennottouse{well above $0.1c$ (synchrotrons, the LHC). Then put $p_\mathrm{v}$ in
    place of $mv$: $r = \dfrac{p_\mathrm{v}}{qB}$. The period grows, so the ions fall out
    of step with a fixed AC.}
  \rearranged{$v = \dfrac{qBr}{m}$ \qquad $B = \dfrac{mv}{qr}$ \qquad
    $f = \dfrac{1}{T} = \dfrac{qB}{2\pi m}$}
  \tensecond{proton in \qty{1.0}{T}: $T = \dfrac{2\pi \times \num{1.673e-27}}{\num{1.602e-19}
    \times 1.0} = \qty{6.6e-8}{s}$, so $f = \qty{1.5e7}{Hz}$.}
  \derive{the magnetic force ($\theta = \ang{90}$) is the centripetal force (both on the
    sheet).
    \begin{align*}
      qvB &= \frac{mv^2}{r} && \text{magnetic force $=$ centripetal force}\\[4pt]
      r &= \frac{mv^2}{qvB} && \text{multiply both sides by $r$, then divide by $qvB$}\\[4pt]
      r &= \frac{mv}{qB} && \text{one $v$ cancels (top and bottom)}\\[4pt]
      T &= \frac{2\pi r}{v} = \frac{2\pi}{v} \times \frac{mv}{qB} = \frac{2\pi m}{qB}
        && \text{from $v = \dfrac{2\pi r}{T}$; the $v$ cancels}\\[4pt]
      f &= \frac{1}{T} = \frac{qB}{2\pi m} && \text{$f = \dfrac{1}{T}$ is on the sheet}
    \end{align*}}
  \trap{use the radius, not the diameter (the largest $r$ is the dee radius). Keep the
    $2\pi$: $f = \dfrac{qB}{2\pi m}$, not $\dfrac{qB}{m}$.}
\end{formulacard}

\begin{workedexample}{A hospital cyclotron}\label{ch5:we:cyclotron}
A hospital needs \qty{11.0}{MeV} protons to make fluorine-18 for PET scans. Its cyclotron
magnet gives \qty{1.2}{T}. Find (a) the AC frequency and (b) the smallest dee radius.
\given{$K = \qty{11.0}{MeV}$, \ $B = \qty{1.2}{T}$, \ $q = \qty{1.602e-19}{C}$, \
  $m = \qty{1.673e-27}{kg}$}
\find{(a) $f$ \ (b) the radius $r$ of the last half circle (the dee must be at least this
  big)}
\formula{$f = \dfrac{1}{T} = \dfrac{qB}{2\pi m}$; \ $K = \tfrac12 mv^2$, so
  $v = \sqrt{\dfrac{2K}{m}}$; \ then $r = \dfrac{mv}{qB}$}
\substitute{\eqlines{f &= \frac{\qty{1.602e-19}{C} \times \qty{1.2}{T}}{2\pi \times
  \qty{1.673e-27}{kg}} = \qty{1.829e7}{Hz}\\
  K &= \num{11.0e6} \times \qty{1.602e-19}{J} = \qty{1.762e-12}{J}
  \quad\text{(eV to J)}}\par
  \eqlines{
  v &= \sqrt{\frac{2 \times \qty{1.762e-12}{J}}{\qty{1.673e-27}{kg}}}
  = \qty{4.590e7}{m.s^{-1}}\\
  r &= \frac{\qty{1.673e-27}{kg} \times \qty{4.590e7}{m.s^{-1}}}{\qty{1.602e-19}{C}
  \times \qty{1.2}{T}} = \qty{0.3994}{m}}}
\answer{(a) \final{$f = \qty{1.8e7}{Hz}$} \ (b) \final{$r = \qty{0.40}{m}$} (2 significant
  figures, like $B$)}
\checkstep{$v \approx 0.15c$, a little above $0.1c$. Relativity makes $r$ only 0.3\%
  bigger here, so $\tfrac12 mv^2$ is still fine to 2 significant figures. The gap voltage
  is not needed: more voltage only means fewer turns.}
\end{workedexample}

Near $c$, the momentum grows much faster than the speed (relativity). So each turn takes
longer, and the ions fall out of step with the AC. A \term{synchrocyclotron} lowers the AC
frequency as the ions speed up. Cyclotrons are small enough to fit in a hospital. There
they make radioactive isotopes for medical scans (for example, fluorine-18 for PET scans)
and proton beams that destroy tumours.

\subsection{The synchrotron and the Large Hadron Collider}

A cyclotron's magnet must cover the whole spiral. Since $r = \dfrac{p_\mathrm{v}}{qB}$,
very high momentum would need a gigantic magnet. A \term{synchrotron} is a ring where the
magnetic field is raised as the particles gain momentum, so the radius stays the same.
The pushes are also timed to match each lap: the ``synchro'' in the name
(\cref{ch5:fig:sync}).

\begin{sidebyside}[0.46]
  \centering
  \begin{tikzpicture}[x=0.82cm, y=0.82cm]
    \draw[hscvel, line width=0.8pt, midarrow] (1.95,0) arc[start angle=0, end angle=359.9, radius=1.95];
    \draw[hscforce, line width=0.8pt, postaction={decorate}, decoration={markings,
      mark=at position 0.3 with {\arrow{Stealth[length=2.2mm, width=1.7mm, reversed]}}}]
      (2.12,0) arc[start angle=0, end angle=359.9, radius=2.12];
    \foreach \a in {30, 60, 120, 150, 210, 240, 300, 330}
      \fill[hscbrand, rotate=\a] (1.87,-0.16) rectangle (2.2,0.16);
    \draw[hscacc, fill=hscacc!25, line width=0.8pt] (-0.3,-2.27) rectangle (0.3,-1.8);
    \draw[hscink, fill=hscseriesC!25, line width=0.7pt] (2.03,0) circle (0.3);
    \draw[hscink, fill=hscseriesC!25, line width=0.7pt] (-2.03,0) circle (0.3);
    % inside each detector the two beams are steered together and cross
    \draw[hscvel, line width=0.8pt] (1.84,-0.17) -- (2.22,0.17);
    \draw[hscforce, line width=0.8pt] (1.84,0.17) -- (2.22,-0.17);
    \draw[hscvel, line width=0.8pt] (-2.22,-0.17) -- (-1.84,0.17);
    \draw[hscforce, line width=0.8pt] (-2.22,0.17) -- (-1.84,-0.17);
    \draw[hscink, fill=hscmuted!20, line width=0.6pt] (-2.4,2.55) rectangle (-1.2,2.8);
    \draw[hscink, line width=0.7pt] (0.25,2.68) circle (0.35);
    \draw[->, hscink, line width=0.6pt] (-1.2,2.68) -- (-0.1,2.68);
    \draw[->, hscink, line width=0.6pt] (0.25,2.33) -- (0.25,2.14);
    \node[font=\scriptsize, anchor=south] at (-1.8,2.8) {linac};
    \node[font=\scriptsize, anchor=west] at (0.62,2.68) {booster ring};
    \node[font=\scriptsize, text=hscbrand] at (0,0.75) {bending magnets};
    \draw[hscbrand!60, line width=0.4pt] (0.3,0.95) -- (0.85,1.45);
    \draw[hscbrand!60, line width=0.4pt] (-0.3,0.95) -- (-0.85,1.45);
    \node[font=\scriptsize, text=hscacc!80!black] at (0,-1.4) {RF cavity};
    \node[font=\scriptsize, align=center] at (0,-0.3) {two beams go\\[-2pt]opposite ways};
    \node[font=\scriptsize, anchor=west] at (2.38,0) {detector};
    \node[font=\scriptsize, anchor=east] at (-2.38,0) {detector};
  \end{tikzpicture}
  \captionof{figure}{A synchrotron collider (not to scale). The two beams are steered
    together inside each detector, where they collide. In the LHC both beams are protons,
    so they travel in two separate pipes with opposite magnetic fields.}\label{ch5:fig:sync}
\nextside
  \begin{itemize}
    \item A linac and smaller rings speed the particles up first.
    \item An \term{RF cavity} (RF = radio frequency) is a metal box with a fast-changing
          electric field. It pushes the particles once per lap.
    \item Bending magnets keep the particles on the ring. $B$ is raised as the momentum
          grows, so $r = \dfrac{p_\mathrm{v}}{qB}$ stays fixed.
    \item Other magnets focus the beam into a thin line.
  \end{itemize}
\end{sidebyside}

A \term{collider} is an accelerator where two beams meet head-on. Think of two equal cars
crashing head-on at the same speed. Their total momentum is zero, so the wreck stops and
all the kinetic energy goes into the crash. When a car hits a parked car, the wreck keeps
moving, so much of the energy stays as motion. In the same way, head-on collisions let
much more of the energy become new particles.

The \term{Large Hadron Collider} (LHC) at CERN, near Geneva, is a \qty{27}{km} ring. Since
2022 each proton in the LHC has \qty{6.8}{TeV} of energy. When two protons meet head-on,
the collision has $6.8 + 6.8 = \qty{13.6}{TeV}$. The LHC's magnets give about \qty{8}{T}.
The protons move only about \qty{3}{m.s^{-1}} slower than light: extra energy adds
momentum, not much speed. \Cref{ch5:tab:compare} compares the three machines.

\begin{table}
  \caption{Three kinds of accelerator.}\label{ch5:tab:compare}
  \centering
  \small
  \begin{tabular}{@{}L{20mm}L{36mm}L{36mm}L{38mm}@{}}
    \toprule
     & \colhead{Linac} & \colhead{Cyclotron} & \colhead{Synchrotron}\\
    \midrule
    Path & straight line & spiral outwards & ring of fixed radius\\
    Speeds up by & electric field in the gaps between tubes & electric field in the gap
      between the dees & electric field in RF cavities, once per lap\\
    Magnetic field & only to focus & constant; bends the path & raised as the momentum
      grows\\
    Limit & length & falls out of step near $c$ & size and magnet strength\\
    Used for & X-ray cancer treatment; first stage & medical isotopes; proton therapy &
      particle physics (LHC); X-ray light sources\\
    \bottomrule
  \end{tabular}
\end{table}

\subsection{What accelerators have shown}

A \term{detector} surrounds the collision point. It tracks and measures what comes out.
Curved tracks in a magnetic field give each particle's charge sign and momentum
($r = \dfrac{p_\mathrm{v}}{qB}$). Neutral particles leave no track. Most new particles
decay at once, so they are found from what they decay into. Each row of
\cref{ch5:tab:tests} is a prediction that an accelerator then confirmed.

\begin{table}
  \caption{Standard Model predictions tested by accelerators.}\label{ch5:tab:tests}
  \centering
  \small
  \begin{tabular}{@{}L{50mm}L{46mm}L{42mm}@{}}
    \toprule
    \colhead{Theory predicts} & \colhead{Accelerator test} & \colhead{Result}\\
    \midrule
    the $\Omega^-$, charge $-1$, and its mass (Gell-Mann's pattern, 1962; later explained
      as sss) & Brookhaven, 1964 & found, as predicted\\
    quarks inside protons (1964) & SLAC linac, deep inelastic scattering, 1968 &
      point-like charges inside\\
    W and Z bosons, about 80 and \qty{91}{GeV}\chFiveperc & CERN proton--antiproton
      collider, 1983 & found at the predicted masses\\
    a third-generation quark (top) & Fermilab (USA), 1995 & found,
      \qty{173}{GeV}\chFiveperc\\
    the Higgs boson (1964) & LHC, CERN, 2012 & found, \qty{125}{GeV}\chFiveperc\\
    \bottomrule
  \end{tabular}
\end{table}

The Higgs idea (1964) came first. The theory used it to predict the masses of the W and
Z. The W and Z were found at those masses in 1983, and the Higgs boson itself only in
2012. Experiments can also lead to new theory: the particle zoo led to the quark model.
Accelerators test Module 7 too: particles never reach $c$, and fast unstable particles
live longer (time dilation).

Particle masses are often given in \si{GeV}\chFiveperc. A mass of \qty{1}{GeV}\chFiveperc\
means the particle's rest energy, $mc^2$, is \qty{1}{GeV}. To change it to kilograms, use
$E = mc^2$ (on the sheet; its formula card is in \cref{ch1:sec:bigbang}).

\begin{workedexample}{How heavy is the Higgs boson?}\label{ch5:we:higgs}
The Higgs boson has a mass of \qty{125}{GeV}\chFiveperc. Find its mass in kilograms. How
many times heavier is it than a proton?
\given{rest energy $E = \qty{125}{GeV} = \qty{125e9}{eV}$, \ $c = \qty{3.00e8}{m.s^{-1}}$,
  \ $m_\mathrm{p} = \qty{1.673e-27}{kg}$}
\find{the mass $m$ in \si{kg}, then $m \div m_\mathrm{p}$}
\formula{$E = mc^2$, so $m = \dfrac{E}{c^2}$ (divide both sides by $c^2$)}
\substitute{\eqlines{E &= \num{125e9} \times \qty{1.602e-19}{J} = \qty{2.0025e-8}{J}
  \quad\text{(eV to J)}\\
  m &= \frac{\qty{2.0025e-8}{J}}{(\qty{3.00e8}{m.s^{-1}})^2}
  = \frac{\qty{2.0025e-8}{J}}{\qty{9.00e16}{m^2.s^{-2}}} = \qty{2.225e-25}{kg}\\
  \frac{m}{m_\mathrm{p}} &= \frac{\qty{2.225e-25}{kg}}{\qty{1.673e-27}{kg}} = 133}}
\answer{\final{$m = \qty{2.23e-25}{kg}$}, about \final{133 times} the mass of a proton.}
\checkstep{each LHC collision has \qty{13.6}{TeV}, about 109 times the Higgs boson's rest
  energy. So the LHC has enough energy to make it \cmark.}
\end{workedexample}

\subsection{Why our understanding is still incomplete}

\begin{itemize}
  \item \textbf{Gravity} is not in the Standard Model, and the graviton has never been found.
  \item \textbf{Dark matter and dark energy:} ordinary matter is only about 5\% of the
        Universe. \term{Dark matter} (unseen matter known only by its gravity) is about
        27\%. Dark energy (an unknown cause of the Universe's speeding-up expansion) is
        about 68\%. No Standard Model particle explains them.
  \item \textbf{Antimatter:} the Big Bang should have made equal amounts of matter and
        antimatter. But today the Universe is almost all matter. No one knows why
        (\cref{ch1:sec:bigbang}).
  \item \textbf{Neutrino mass:} the first Standard Model had massless neutrinos.
        Experiments since 1998 show they have a tiny mass.
  \item \textbf{Unanswered questions:} why are there exactly three generations? The theory
        cannot predict about 19 of its own numbers, such as particle masses. They must be
        measured.
\end{itemize}

\begin{mistakes}
\begin{itemize}
  \item \xmark\ ``The magnets speed the particles up.'' \cmark\ Electric fields do; magnets
        only bend and focus.
  \item \xmark\ ``Drift tubes get longer to keep the force the same.'' \cmark\ They keep
        the time in each tube equal to half an AC cycle.
  \item \xmark\ ``A cyclotron's frequency must change as the ions speed up.''
        \cmark\ $T = \dfrac{2\pi m}{qB}$ does not depend on speed, so one fixed frequency
        works. The period only starts to grow when the ions get close to $c$.
  \item \xmark\ ``Making new particles breaks conservation of mass.'' \cmark\ Mass--energy
        is conserved: kinetic energy becomes mass ($E = mc^2$).
\end{itemize}
\end{mistakes}

\begin{examtip}
``Explain how a cyclotron works'' (3--4 marks): the gap's electric field speeds the
protons up ($W = qV$). The magnetic field bends them; faster means a bigger radius. The
period does not depend on speed, so a fixed AC frequency stays in step. Another common
question asks you to analyse how theory and experiment together built the Standard Model
(6 marks in the 2025 HSC). Give 2--3 \emph{prediction $\to$ test} pairs
(\cref{ch5:tab:tests}). Add one case where experiment came first (particle zoo $\to$
quarks). Finish with the implication: theory and experiment drive each other, and new
results (neutrino mass, dark matter) show the model is still incomplete.
\end{examtip}

\begin{quickcheck}
  \item An alpha particle (charge $2e$) is accelerated through \qty{50}{kV}. Find the energy
        it gains in \si{keV} and in \si{J}.
  \item A cyclotron for protons uses a \qty{0.80}{T} field. Find the AC frequency. What
        happens to it if $B$ is doubled?
  \item The magnetic field in a cyclotron does no work. Explain why the cyclotron still
        needs it.
  \item An electron and a positron collide head-on and make a tau and an antitau. Each has
        a mass of \qty{3.167e-27}{kg}. Use $E = mc^2$ (as in \cref{ch1:we:proton}) to
        find the least total energy needed, in \si{GeV}.
  \item Give three reasons why our understanding of matter is still incomplete.
\end{quickcheck}

%% ==================================================================
\section{Key experiments}\label{ch5:sec:experiments}
\covers{M8.IQ5.01, M8.IQ5.03}

\begin{experiment}{Deep inelastic scattering at SLAC (1967--1973)}
\aim{To find out whether the proton has parts inside it.}
\apparatus{the \qty{3.2}{km} SLAC electron linac. A target of liquid hydrogen
  (protons), or of liquid deuterium for neutrons (deuterium is heavy hydrogen: one proton
  and one neutron). Magnetic spectrometers that measure each scattered electron's energy
  and angle. A beam dump: a block that safely stops the unscattered beam.
  \begin{figure}
    \centering
    \begin{tikzpicture}[x=1cm, y=1cm]
      \draw[draw=hscink, fill=hscmuted!18, line width=0.7pt] (0,-0.25) rectangle (4.0,0.25);
      \node[font=\scriptsize] at (2.0,0) {electron linac (\qty{3.2}{km})};
      \draw[velocity] (4.0,0) -- (5.55,0);
      \node[font=\scriptsize, text=hscvel, anchor=south] at (4.75,0.05) {electrons};
      \draw[draw=hscink, fill=hscseriesA!20, line width=0.7pt] (5.8,0) circle (0.25);
      \node[font=\scriptsize, align=center, anchor=north] at (5.8,-0.32) {liquid hydrogen\\[-2pt]target};
      \draw[guide] (6.05,0) -- (9.3,0);
      \draw[draw=hscink, fill=hscink!70, line width=0.7pt] (9.3,-0.3) rectangle (9.7,0.3);
      \node[font=\scriptsize, anchor=north] at (9.5,-0.35) {beam dump};
      \draw[velocity] (5.8,0) -- (7.75,1.13);
      \draw[draw=hscink, fill=hscbrand!12, line width=0.7pt, rotate around={30:(8.35,1.48)}]
        (7.75,1.18) rectangle (8.95,1.78);
      \node[font=\scriptsize, align=center, anchor=west] at (9.05,1.75)
        {magnetic spectrometer:\\[-2pt]measures energy\\[-2pt]and angle};
      \draw[hscink, line width=0.5pt] (7.0,0) arc[start angle=0, end angle=30, radius=1.2];
      \node[font=\small] at (7.25,0.33) {$\theta$};
      \draw[->, hscmuted, line width=0.5pt, dashed] (8.7,0.55) arc[start angle=8, end angle=50, radius=2.9];
      \node[font=\scriptsize, text=hscmuted, anchor=east] at (7.5,2.42) {swings to new angles};
    \end{tikzpicture}
    \caption{Deep inelastic scattering at SLAC (top view, not to scale).}\label{ch5:fig:slac}
  \end{figure}}
\method{fire electrons (up to about \qty{20}{GeV}) at the target. At each angle $\theta$,
  count the scattered electrons and measure the energy each lost. Repeat for many angles
  and energies, and compare with the prediction for a smooth proton.}
\results{many electrons scattered at large angles and lost a lot of energy. A smooth
  proton would give far fewer of these. The electrons acted as if they hit tiny
  point-like charges.}
\conclusion{protons and neutrons contain point-like charged parts: quarks.}
\significance{the first direct evidence of point-like parts (quarks) inside protons and
  neutrons. It turned the quark model from a pattern into real particles (Nobel Prize
  1990).}
\validity{\begin{itemize}
  \item \textbf{Validity:} electrons have no parts and ignore the strong force. Their
        wavelength (about \qty{e-16}{m}; \qty{6e-17}{m} at \qty{20}{GeV}) is much smaller
        than a proton.
  \item \textbf{Reliability:} very many electrons counted at many angles and energies;
        later experiments at CERN and DESY (Germany) agreed.
  \item \textbf{Accuracy:} precise magnetic spectrometers, with small corrections for
        energy lost before and after scattering.
\end{itemize}}
\end{experiment}

\begin{experiment}{Discovery of the Higgs boson at the LHC (2012)}
\aim{To find the Higgs boson predicted in 1964, and measure its mass.}
\apparatus{the LHC (\cref{ch5:fig:sync}) and two huge detectors, ATLAS and CMS
  (\cref{ch5:fig:detector}). Each detector is built in layers around the beam pipe.
  \begin{figure}
    \centering
    \begin{tikzpicture}[x=0.95cm, y=0.95cm]
      \fill[hscmuted!12] (0,0) circle (2.2);
      \draw[hscink, line width=0.6pt] (0,0) circle (2.2);
      \fill[white] (0,0) circle (1.75);
      \draw[hscink, line width=0.6pt] (0,0) circle (1.75);
      \fill[hscseriesB!16] (0,0) circle (1.5);
      \draw[hscink, line width=0.6pt] (0,0) circle (1.5);
      \fill[white] (0,0) circle (0.95);
      \fill[hscseriesA!12] (0,0) circle (0.95);
      \draw[hscink, line width=0.6pt] (0,0) circle (0.95);
      \draw[hscink, fill=white, line width=0.6pt] (0,0) circle (0.16);
      \fill[hscforce] (0,0) circle (1.4pt);
      % two photons: no track in the tracker, energy dumped in the calorimeter
      \draw[hscray, line width=1pt, decorate, decoration={snake, amplitude=0.7mm,
        segment length=2mm}] (0,0) -- (120:1.3);
      \draw[hscray, line width=1pt, decorate, decoration={snake, amplitude=0.7mm,
        segment length=2mm}] (0,0) -- (-60:1.3);
      \fill[hscray] (120:1.3) circle (2.4pt);
      \fill[hscray] (-60:1.3) circle (2.4pt);
      % a charged particle: curved track in the tracker
      \draw[hscvel, line width=0.8pt] (0,0) arc[start angle=60, end angle=110, radius=1.0];
      \foreach \chFiveX/\chFiveR/\chFiveT/\chFiveY in {1.6/1.25/{muon detectors}/1.6,
          1.3/0.45/{calorimeter (measures energy)}/0.75,
          0.55/0.15/{tracker (curved tracks)}/-0.1,
          0.1/-0.1/{beam pipe (beams go into the page)}/-0.95}
        {\draw[hscmuted, line width=0.4pt] (\chFiveX,\chFiveR) -- (2.55,\chFiveY);
         \node[font=\scriptsize, anchor=west] at (2.6,\chFiveY) {\chFiveT};}
      \node[font=\scriptsize, text=hscray!85!black, anchor=north] at (0,-2.25)
        {wavy lines: two photons from a Higgs boson};
    \end{tikzpicture}
    \caption{Cross-section of a detector such as ATLAS or CMS (sketch, not to scale).}\label{ch5:fig:detector}
  \end{figure}}
\method{collide protons (\qtyrange{7}{8}{TeV} in total) hundreds of millions of times a
  second. Keep the \term{events} (single collisions and what comes out of them) that have
  the products a Higgs should give: two photons, or four leptons. Use their energies and
  momenta to find the mass of the particle that decayed into them. Plot the number of
  events against this mass.}
\results{a small bump at about \qty{125}{GeV}\chFiveperc\ above a smooth background
  (\cref{ch5:fig:bump}), seen by both ATLAS and CMS.
  \begin{figure}
    \centering
    \begin{tikzpicture}
      \begin{axis}[hsc, width=0.6\linewidth, height=0.3\linewidth,
          xmin=105, xmax=160, ymin=0, ymax=1100, ytick=\empty,
          xlabel={mass of the photon pair (\si{GeV}\chFiveperc)},
          ylabel={number of events}]
        \addplot[hscseriesA, line width=1.2pt, domain=105:160, samples=160]
          {1000*exp(-(x-105)/30) + 200*exp(-(x-125)^2/(2*2.2^2))};
        \addplot[hscseriesB, dashed, line width=1pt, domain=105:160, samples=80]
          {1000*exp(-(x-105)/30)};
        \draw[->, hscseriesA, line width=0.6pt] (axis cs:131,880) -- (axis cs:126,725);
        \node[font=\scriptsize, text=hscseriesA, anchor=south west] at (axis cs:131,860)
          {bump at \qty{125}{GeV}\chFiveperc: the Higgs};
        \draw[->, hscseriesB!80!black, line width=0.6pt] (axis cs:122,250) -- (axis cs:124.5,500);
        \node[font=\scriptsize, text=hscseriesB!80!black, anchor=north] at (axis cs:122,250)
          {background (other events)};
      \end{axis}
    \end{tikzpicture}
    \caption{Sketch of the kind of graph shown in 2012 (not real data).}\label{ch5:fig:bump}
  \end{figure}}
\conclusion{a new boson of about \qty{125}{GeV}\chFiveperc\ exists: the Higgs boson.}
\significance{the last missing Standard Model particle. The theorists Englert and Higgs,
  who predicted it, won the Nobel Prize in 2013.}
\validity{\begin{itemize}
  \item \textbf{Validity:} it looked for the exact decay products the theory predicts (two
        photons; four leptons). It compared the counts with the background expected with
        no Higgs.
  \item \textbf{Reliability:} two separate detectors and teams found the same mass. The
        chance of a random bump had to be below 1 in 3.5 million (``5 sigma'').
  \item \textbf{Accuracy:} calibrated with known particles (the Z). In 2012 the mass was
        known to about 0.5\%. Today it is known to about 0.1\%
        (\qty{125.2}{GeV}\chFiveperc).
\end{itemize}}
\end{experiment}
